\documentclass[11pt,a4paper]{article}
\usepackage[margin=2.4cm]{geometry}
\usepackage{amsmath,amssymb,booktabs,graphicx,hyperref,xcolor}
\usepackage[T1]{fontenc}
\hypersetup{colorlinks=true,linkcolor=blue!50!black,citecolor=blue!50!black,
            urlcolor=blue!50!black}
\newcommand{\Szero}{\ensuremath{\Sigma^{0}}}
\newcommand{\Lam}{\ensuremath{\Lambda}}
\newcommand{\uboone}{MicroBooNE}
\newcommand{\numubar}{\ensuremath{\bar{\nu}_\mu}}

\title{\Szero\ Production in the NuMI Beam at \uboone:\\ Analysis Note}
\author{J.~Nowak, N.~Patel}
\date{\today}

\begin{document}
\maketitle

\begin{abstract}
Working note for the \uboone\ \Szero\ production analysis in the NuMI beam. The
selection reconstructs the $\Szero \to \Lam\gamma$, $\Lam \to p\pi^-$ chain in
$\numubar$ charged-current interactions and is finished with a cross-validated
gradient BDT. This note records the current state, the planned port into the
\texttt{xsec\_analyzer} framework, and the three pieces of work that follow:
a Poisson-thrown fake-data sample, extension from RHC-only to the combined
FHC$+$RHC exposure, and the selection work needed to lift a purity that currently
sits near $0.1\%$.
\end{abstract}

\tableofcontents

% =============================================================================
\section{Status and scope}
% =============================================================================

The analysis measures, or more realistically sets a limit on, \Szero\ production
in $\numubar$ charged-current interactions. The signature is
\begin{equation}
  \numubar + \mathrm{Ar} \to \mu^{+} + \Szero + X, \qquad
  \Szero \to \Lam\gamma, \qquad \Lam \to p\,\pi^{-},
\end{equation}
so the reconstructed final state is a muon, a proton, a pion and a photon.

The existing selection (\texttt{Sigma/pelee/SigmaSelection.C}) runs on standard
PeLEE ntuples and reaches, at the operating point $\mathrm{BDTG}>-0.90$ scaled to
the full RHC data exposure of $1.1526\times10^{21}$~POT:

\begin{table}[h]\centering
\caption{Current selection performance, RHC only.}
\begin{tabular}{lrrrr}
\toprule
Stage & Raw signal & Efficiency [\%] & Purity [\%] & $S/\sqrt{S+B}$ \\
\midrule
Preselection (NTracksCut) & 650 & 15.7 & 0.040 & 0.038 \\
Cut-based                 & 497 & 12.0 & 0.077 & 0.045 \\
BDTG (current final)      & 502 & 12.1 & 0.108 & 0.054 \\
\bottomrule
\end{tabular}
\end{table}

\noindent Efficiency is respectable; purity is not. A purity of $0.108\%$ means
roughly one signal event per thousand selected, and $\sim\!2640$ background events
under $\sim\!2.9$ signal events at full RHC POT. \textbf{Improving the selection is
therefore the priority}, ahead of any normalisation or extraction work: at this
purity no downstream treatment can rescue the measurement.

\paragraph{What this analysis produces.} Not a differential cross section. At
$0.1\%$ purity and $12\%$ efficiency the deliverable is a cross-section
\emph{upper limit}. This matters for the framework port (Sec.~\ref{sec:port}),
because \texttt{xsec\_analyzer} is built end-to-end for differential extraction and
contains no limit machinery.

% =============================================================================
\section{Framework port}
\label{sec:port}
% =============================================================================

The selection is being ported into \texttt{xsec\_analyzer} as a
\texttt{SelectionBase} subclass, alongside the CC$1\mu1\pi$ selections. The port
buys the framework's systematic machinery -- flux (PPFX multisims),
cross-section (GENIE multisims and unisims), detector variations, hadron
reinteraction -- which is exactly the covariance a limit needs for its background
prediction, and which would otherwise have to be built from scratch.

\paragraph{What is ported.} Selection, cut flow, truth categorisation, reco and
true observables, and the systematic universes at a \emph{single} analysis bin.

\paragraph{What is not.} Limit setting, which stays in \texttt{CrossSection.C}
reading the framework covariance. An unfolded \Szero\ spectrum would be
meaningless at this purity, so no multi-bin scheme is planned.

\paragraph{Framework gaps.} Three, all small. Nine branches the selection uses are
present in the ntuples but unbound in \texttt{Branches.hh}
(\texttt{trk\_pid\_chipi\_v}, \texttt{shr\_energy\_y\_v},
\texttt{shr\_p\{x,y,z\}\_v}, \texttt{pi0\_energy\{1,2\}\_Y},
\texttt{pi0\_gammadot}, \texttt{shr\_moliere\_avg\_v}); the fiducial volume is
looser than the CC$\pi$ one and must be set per selection rather than inherited;
and the factory needs a registration branch. Details in
\texttt{SIGMA0\_FRAMEWORK\_PLAN.md}.

% =============================================================================
\section{Selection}
% =============================================================================

\subsection{Signal definition}
$\texttt{nu\_pdg} = -14$, $\texttt{ccnc} = 0$, a \Szero\ (PDG 3212) in the
final-state list \texttt{mc\_pdg}, and the true vertex inside the fiducial volume.
Photons are not stored in \texttt{mc\_pdg}, so the decay is tagged by the \Szero\
itself rather than by its daughters.

\subsection{Fiducial volume}
$3.00 < x < 253.35$, $-112.53 < y < 114.47$, $3.1 < z < 1033.0$~cm, excluding the
dead region $675.1 < z < 775.1$~cm. This is deliberately looser than the
CC$1\mu1\pi$ box ($10$--$246$, $\pm101$, $10$--$986$): the signal is rare enough
that acceptance is worth more than vertex quality.

\subsection{Candidate reconstruction}
\begin{itemize}
  \item \textbf{Muon}: longest track with $\texttt{trk\_score}>0.8$, start within
        $20$~cm of the vertex, length $>10$~cm, and LLR PID $>0.7$.
  \item \textbf{Proton and pion}: assigned jointly by the combined discriminant
        \begin{equation}
          D = \frac{\mathrm{LLR}}{0.591}
              + \frac{\chi^2_p - \chi^2_\pi}{124.5},
          \label{eq:Ddisc}
        \end{equation}
        with proton $=\arg\min D$ and pion $=\arg\max D$. This separates true
        protons from true pions with AUC $0.873$, against $0.838$ for $\chi^2_p$
        alone and $0.852$ for LLR alone. The $\chi^2$ term is dropped per track
        when undefined ($\sim\!30\%$ of tracks).
  \item \textbf{Photon}: the two non-track PFPs closest to the vertex; the closest
        is taken as the \Szero\ decay photon.
\end{itemize}

\subsection{Cut chain}
Nine cumulative stages: \texttt{AllEvents}, \texttt{FiducialV},
\texttt{CosmicVeto} ($\texttt{topological\_score}>0.20$), \texttt{NShowers}
($0<n_{\mathrm{shr}}<4$), \texttt{MuonCandidate}, \texttt{ProtonCandidate},
\texttt{PionCandidate}, \texttt{NTracksCut} ($n_{\mathrm{trk}}<5$), \texttt{BDTG}.

\subsection{Derived observables}
\Lam\ mass from $(p,\pi)$; \Szero\ mass from $(\Lam,\gamma)$ using the per-PFP
shower energy \texttt{shr\_energy\_y\_v}; \Lam\ and photon vertex-pointing angles
and impact parameters; \Lam\ decay distance; and a $\pi^0$ di-photon mass built
from the framework $\pi^0$ reconstruction for background rejection.

\subsection{BDT}
Gradient BDT over 49 inputs, evaluated with two-fold cross-validation: folds are
trained on even and odd event numbers and each is applied only to the events it did
not train on. Per-fold ROC $\simeq0.81$. The cross-validation is not optional --
applying a fold to its own training events would inflate the ROC and the apparent
significance.

% =============================================================================
\section{Why the purity is low, and what to do about it}
\label{sec:purity}
% =============================================================================

This is the priority. Two diagnostics point at the same place.

\paragraph{The surviving background is $\pi^0$ events.}
\begin{table}[h]\centering
\caption{Surviving background by interaction mode at the final selection.}
\begin{tabular}{lrr}
\toprule
Mode & Scaled yield & \% of background \\
\midrule
Resonant (RES)       & 1381.8 & 52.3 \\
Deep inelastic (DIS) &  635.3 & 24.0 \\
Quasi-elastic (QE)   &  363.8 & 13.8 \\
Meson exchange (MEC) &  177.1 &  6.7 \\
Cosmic               &   74.6 &  2.8 \\
Coherent (COH)       &    8.9 &  0.3 \\
\midrule
Total                & 2643.1 & 100 \\
\bottomrule
\end{tabular}
\end{table}

\noindent RES and DIS together are $76\%$. Both produce $\pi^0 \to \gamma\gamma$,
which supplies a fake \Szero\ photon, plus charged hadrons that fake the proton and
pion tracks. The background is not diffuse: it is one physics process.

\paragraph{The candidate assignment is the weak link.} Truth-matched purities of
the individual candidates in background events:

\begin{center}
\begin{tabular}{lrl}
\toprule
Candidate & Correct [\%] & Dominant impostor \\
\midrule
Muon   & 81 & --- \\
Proton & 58 & --- \\
Pion   & 31 & another proton (28\%) \\
Photon & 29 & a proton (18\%) \\
\bottomrule
\end{tabular}
\end{center}

\noindent The muon is reliable and the proton is acceptable, but \textbf{the pion
and photon candidates are wrong roughly $70\%$ of the time}. Since both the \Lam\
and \Szero\ masses are built from them, the two most powerful BDT inputs are being
computed from misidentified objects in most background events -- and, importantly,
we do not know how often in signal events, because the NoCR signal sample's
non-muon truth matching is unreliable.

\subsection{The photon candidate: what the data says}
\label{sec:photonstudy}

This was investigated first, and the result contradicts the obvious expectation, so
it is recorded in full.

\paragraph{Composition.} Truth-matching the photon candidate in background events
(69557 preselected events) gives:

\begin{center}
\begin{tabular}{lr}
\toprule
Truth match & Fraction [\%] \\
\midrule
photon              & 30.0 \\
unmatched (PDG 0)   & 26.7 \\
proton              & 15.2 \\
charged pion        & 16.3 \\
muon                &  9.8 \\
electron            &  1.4 \\
\bottomrule
\end{tabular}
\end{center}

\noindent Charged tracks reconstructed as showers -- proton, pion, muon together --
account for $41\%$, and a further $27\%$ has no truth match at all.

\paragraph{The signal photon is soft, which inverts the strategy.} The
$\Szero$--$\Lam$ mass splitting is only $77$~MeV, so
\begin{equation}
  E_\gamma = \frac{m_{\Szero}^2 - m_{\Lam}^2}{2 m_{\Szero}} = 74.5~\mathrm{MeV}
\end{equation}
in the \Szero\ rest frame. The measured mean energy of the photon candidate in
signal events is $70$~MeV, which matches that prediction and is good evidence that
the selection is finding the right object. The background photon candidate averages
$140$~MeV, because the $\pi^0$ photons that dominate it are far harder.

\textbf{Requiring shower-like quality would therefore remove signal
preferentially}, which is the opposite of what an earlier draft of this note
proposed. A $74$~MeV photon is near the reconstruction threshold: it makes few
hits and deposits little energy. Measured against true photons in background
events -- which are $\pi^0$ photons and much harder -- the signal photon looks
\emph{less} shower-like, not more.

\paragraph{Cutting in the correct direction helps, but only before the BDT.}
Scanning cuts in the soft direction at preselection:

\begin{center}
\begin{tabular}{lrrr}
\toprule
Cut & $\epsilon_S$ [\%] & Background rejected [\%] & FoM gain \\
\midrule
$E_\gamma < 120$~MeV          & 94.5 & 25.4 & $\times1.092$ \\
$N_{\mathrm{hits}}^{Y} < 102$ & 95.7 & 21.5 & $\times1.079$ \\
$E_\gamma \in [5,100]$~MeV    & 70.6 & 55.2 & $\times1.052$ \\
\bottomrule
\end{tabular}
\end{center}

\noindent But \texttt{Shower1Energy} is already one of the 49 BDT inputs, so the
question is whether the cut adds anything the BDT has not already taken. Applied on
top of the BDT operating point it does not:

\begin{center}
\begin{tabular}{lrrr}
\toprule
Selection & $S$ & $B$ & FoM \\
\midrule
BDT only                            & 502 & 20762 & 3.443 \\
BDT $+$ $E_\gamma<120$~MeV          & 482 & 18663 & 3.484 \quad($\times1.012$) \\
BDT $+$ $E_\gamma<100$~MeV          & 456 & 17706 & 3.384 \quad($\times0.983$) \\
\bottomrule
\end{tabular}
\end{center}

\paragraph{Conclusion.} The photon candidate cannot be improved by cutting on
variables the BDT already has -- it is using them close to optimally. An explicit
$\pi^0$ mass veto is likewise ineffective, rejecting only $9\%$ of background,
because it requires \emph{two} reconstructed showers and most background events
supply only one. Improvement must come from information the BDT does not yet have.

\paragraph{The one genuinely unused handle.} \texttt{shr\_moliere\_avg\_v}, the
shower Moli\`ere angle, is in the ntuples and in the plotting variable list but is
\emph{not} among the 49 BDT inputs. It measures lateral spread, which is
orthogonal to energy and hit count and is the classic discriminant between an
electromagnetic shower and a track stub. Adding it requires re-running the
selection to dump it into the training tree, then retraining both folds. That is
the concrete next step for the photon candidate, and it should be judged by the
FoM gain on top of the BDT, not at preselection.

\subsection{Methods worth adopting from the single-photon analyses}
\label{sec:glee}

The gLEE single-photon search and the NC$\pi^0$ technote solve a problem very close
to ours: isolate a single photon against a background of $\pi^0$ events in which one
decay photon was missed. Four of their methods transfer directly.

\paragraph{1. Second-Shower Veto -- look for the photon that was \emph{not}
reconstructed.} This is the most important one. Our $\pi^0$ veto requires
\emph{both} decay photons to be reconstructed, which is why it fires so rarely and
rejects only $9\%$ of background (Sec.~\ref{sec:photonstudy}). gLEE instead trains a
dedicated Second-Shower Veto (SSV) BDT on evidence that a second photon exists even
when Pandora never built a shower from it. Their SSV inputs include:
\begin{itemize}
  \item the \textbf{fraction of unassociated hits in the slice} -- hits in the
        neutrino slice belonging to no reconstructed object, which is exactly what a
        missed low-energy photon leaves behind;
  \item the \textbf{number of unassociated hits} (used logarithmically);
  \item a \textbf{minimum conversion parameter}, and the implied invariant mass of
        the candidate photon paired with any second-shower candidate.
\end{itemize}
Applied to \Szero, this attacks the dominant background where it actually lives: our
RES and DIS background is $76\%$ of the total and is $\pi^0$-driven, and the failure
mode is precisely a missed second photon. This should be tried before any further
work on the existing $\pi^0$ mass variable.

The ingredients are already in our ntuples, checked directly on the signal file:
\texttt{slnhits} and \texttt{slnunhits} (total and unassociated hits in the slice,
whose ratio is gLEE's fraction variable), \texttt{slclustfrac},
\texttt{secondshower\_Y\_nhit} and \texttt{secondshower\_Y\_dot} from the PeLEE
second-shower tagger, and \texttt{shrsubclusters0/1/2} (sub-cluster counts, which the
CC$1\mu1\pi$ BDT already uses). None of these is currently read by
\texttt{SigmaSelection.C}. Adding them is the same mechanical exercise as the
Moli\`ere angle -- bind the branch, add to both variable lists, retrain -- but on
variables with a far better physics motivation.

\paragraph{2. One BDT per background class, not one BDT for everything.} gLEE runs
\emph{five} BDTs for its $1\gamma1p$ channel -- Cosmic, $\nu_e$, BNB-Other, NC$\pi^0$
and SSV -- each trained against one background, with a separate cut on each score.
We use a single BDT over 49 inputs trained against all backgrounds at once. Given
that our background decomposes cleanly (RES $52\%$, DIS $24\%$, QE $14\%$, MEC
$7\%$, cosmic $3\%$), a dedicated $\pi^0$-rejection BDT and a separate cosmic BDT
would let each learn its own discriminating structure instead of averaging over
processes that look nothing alike. This is an architectural change and the single
largest departure from the current design.

\paragraph{3. Conversion distance as a first-class variable.} gLEE uses the shower
conversion distance throughout, in both linear and logarithmic form, as one of its
most-shown discriminants. We have \texttt{Shower1VtxDist}, which is the distance
from the vertex to the shower start and is closely related, but it is one input among
49 rather than a variable the selection is built around. A real photon converts a few
centimetres from the vertex; a proton track begins at it. Worth examining separately
and in log form, since gLEE's plots show the discrimination lives at small distances
where a linear axis compresses it.

\paragraph{4. Constrain the $\pi^0$ rate from data.} The NC$\pi^0$ technote fits the
$\pi^0$ normalisation to data and finds GENIE over-predicts non-coherent NC$\pi^0$ by
$25\%$: the flat fit gives $N_{\mathrm{non-coh}} = 0.75$ ($\chi^2/\mathrm{ndf} =
1.96/4$), and gLEE separately uses a flat $0.69$ or a momentum-dependent
$0.92 - 0.80\,|\vec{p}_{\mathrm{true}}|$ scaling. Since our background prediction is
dominated by the same physics, an unconstrained $\pi^0$ normalisation is both a bias
and a large systematic. A sideband constraint would reduce both. Note that gLEE
deliberately does \emph{not} apply the scaling in its final fit, relying on the
covariance matrix instead -- that is the more defensible route for us too, and it
argues for making sure the $\pi^0$ normalisation uncertainty is properly represented
in the covariance rather than for rescaling the central value.

\paragraph{What does not transfer.} gLEE's signal is a single photon with no muon,
so its $\nu_e$ and BNB-Other BDTs target backgrounds we do not have; our muon
candidate is already $81\%$ pure and is not where we are losing. Their analysis is
also BNB, so the flux and interaction mix differ, and any borrowed \emph{numerical}
threshold would have to be re-derived.

\subsection{Plan: dedicated BDTs in place of one large BDT}
\label{sec:bdtsplit}

\paragraph{Why the single BDT is the wrong shape.} One BDT over 49 inputs is trained
against a background that is really four unrelated physics processes -- $\pi^0$-producing
RES ($52\%$) and DIS ($24\%$), quasi-elastic ($14\%$) and MEC ($7\%$) events with no
$\pi^0$ at all, and cosmics ($3\%$). A single discriminant must find one surface
separating signal from \emph{all} of them at once, so it learns whatever those
backgrounds have in common and averages away the structure specific to each. gLEE
runs five, and their signal is rarer than ours.

The evidence that this is the binding constraint rather than variable quality is that
adding variables barely moves the answer: the Moli\`ere angle changed the figure of
merit by $\times1.013$, and the ROC plateaus regardless of background size or added
inputs. A better-shaped discriminant is a different lever from a better variable.

\paragraph{Proposed structure.} Three BDTs, cut independently, replacing the single
BDTG:

\begin{center}
\begin{tabular}{llll}
\toprule
BDT & Trained against & Fraction of bkg & Discriminating on \\
\midrule
$\pi^0$-rejection & RES $+$ DIS with a true $\pi^0$ & $\sim\!76\%$ &
  SSV block, $\pi^0$ mass, shower shape \\
Non-$\pi^0$ $\nu$ & QE $+$ MEC $+$ COH, no $\pi^0$ & $\sim\!21\%$ &
  track multiplicity, PID, $\Lam$ mass \\
Cosmic & cosmic-tagged events & $\sim\!3\%$ &
  topological score, containment, $\Lam$ pointing \\
\bottomrule
\end{tabular}
\end{center}

\noindent Each is trained signal-versus-one-class, keeping the same two-fold
cross-validation by event-number parity. The truth labels needed to define the
classes already exist in the training dump: \texttt{Interaction},
\texttt{HasPi0}, \texttt{CosmicMu} and \texttt{Ccnc} are written to
\texttt{BDTInput.root} as diagnostic branches, so the split costs no new
reconstruction, only a change to how \texttt{TrainBDT.C} builds its background tree.

\paragraph{How to combine them.} Two options, and the choice should be made on
measurement rather than taste:
\begin{itemize}
  \item \textbf{Independent cuts}, as gLEE does -- one threshold per score, tuned
        jointly. Transparent, and each cut can be validated in its own sideband.
  \item \textbf{A second-stage combination} -- feed the three scores into a small
        final discriminant. Usually stronger, but harder to validate and easy to
        overtrain on 650 signal events.
\end{itemize}
Given our signal statistics, independent cuts are the safer starting point.

\paragraph{The statistics problem this creates, and the answer to it.} Splitting the
background does not reduce the signal available for training: all $650$ signal events
are used in every BDT. It does reduce the \emph{background} per training, but from
$69557$ preselected events, so even the cosmic class retains thousands. The real
constraint remains the $650$ signal events, which is why the incoming cosmic-overlay
signal samples matter more than any of this.

\paragraph{Order of work.} Build the $\pi^0$-rejection BDT first and evaluate it
against the current single BDT at matched signal efficiency. If it alone does not
improve the figure of merit, the split is not the answer and the effort should move
to signal statistics instead. Only if it does should the other two follow.

\paragraph{Interaction with the SSV variables.} The SSV block belongs naturally in
the $\pi^0$-rejection BDT, where its targets are. Adding it to the single BDT first
-- as has now been done -- is worth doing anyway, because it measures how much of the
gain comes from the variables and how much from the restructuring: if the SSV
variables give little in the single BDT but a lot in a dedicated $\pi^0$ BDT, that is
direct evidence for the split.

\subsection{Proposed selection work, in order}

\begin{enumerate}
  \item \textbf{Second-Shower Veto variables} (Sec.~\ref{sec:glee}): unassociated-hit
        fraction and count in the slice. This targets $76\%$ of the background at its
        actual failure mode and is expected to dominate everything else on this list.
  \item \textbf{Split the BDT by background class} (Sec.~\ref{sec:glee}), starting
        with a dedicated $\pi^0$-rejection BDT.
  \item \textbf{Add \texttt{shr\_moliere\_avg\_v} to the BDT.} \emph{Done, and it
        does almost nothing}, as its separation predicted: $503$ signal against $502$,
        purity $0.111\%$ against $0.108\%$, figure of merit $0.0544$ against $0.0537$,
        a gain of $\times1.013$. Recorded because it closes off the last unused
        photon-shape variable and supports the argument of
        Sec.~\ref{sec:bdtsplit} that variable quality is no longer the binding
        constraint.
  \item \textbf{Improve the pion candidate ($31\%$ correct).} The joint
        discriminant of Eq.~\ref{eq:Ddisc} is already a real improvement over
        $\chi^2_p$ alone, but the failure mode is specific: a second proton is
        taken as the pion $28\%$ of the time. Adding the sign-sensitive
        information -- track length against the proton range curve, or the
        $\Lam$-mass constraint itself as part of the assignment rather than as a
        consequence of it -- would target that directly.
  \item \textbf{Reconsider the \Lam\ decay topology.} \Lam\ has $c\tau = 7.89$~cm,
        so a genuine $\Lam \to p\pi^-$ has a displaced vertex. The selection
        computes \texttt{LambdaDecayDist} and the pointing angle but does not cut
        on them. A displaced-vertex requirement is close to background-free in
        principle and is the most distinctive feature of the signal.
\end{enumerate}

\noindent The analysis is also \textbf{signal-statistics limited}: 650 raw signal
events at preselection, with a BDT ROC that plateaus regardless of background size
or added variables. More signal MC is the largest single lever, and no amount of
selection work substitutes for it.

% =============================================================================
\section{Reconstruction study on HyperonProduction ntuples}
\label{sec:recostudy}
% =============================================================================

\subsection{Why the PeLEE ntuples cannot answer these questions}

Every question about \emph{which particle a candidate actually is} was
unanswerable on the PeLEE ntuples. Their truth record stops at the primary
vertex: \texttt{mc\_pdg} lists the \Szero\ and the muon, but never the
$\Szero \to \Lam\gamma \to p\pi^-\gamma$ chain. Measured on the signal sample,
\texttt{nproton} averages $0.06$ and is zero in $94.4\%$ of events,
\texttt{npion} is zero in $100\%$, PDG~22 never appears, and
\emph{zero of 17721 events} contain a backtracked $\pi^-$ PFP. The
per-candidate truth columns are correspondingly meaningless: the pion candidate
matches a true $\pi^-$ in $0.0\%$ of events, $60\%$ carry no truth match at all
and $36.5\%$ leak the muon's truth. This is a property of the ntuple, not a bug
in the selection.

The HyperonProduction (``hypana'') analyser writes the full decay chain
(\texttt{mc\_sigmazero\_lambda\_end\_*} is the true \Lam\ decay vertex,
\texttt{mc\_decay\_*} the daughters) together with per-PFP truth matching
(\texttt{pfp\_true\_pdg}, \texttt{pfp\_purity}, \texttt{pfp\_completeness}).
The studies below are all performed there. \texttt{HypInput.h} adapts the hypana
branches to the PeLEE names the selection already uses; the two formats index
every \texttt{trk\_*}/\texttt{shr\_*}/\texttt{pfp\_*} vector per PFP with equal
lengths, so the port is a branch remap with a double$\rightarrow$float copy
rather than a structural rewrite.

Two caveats apply throughout. The hyperon samples are generated with an
\emph{enhanced} hyperon-production cross section, so they are valid for
efficiencies and ratios but not for absolute yields. And a full hypana-only
analysis is not yet possible: hypana output exists only for hyperon samples,
there is none for the generic RHC overlay or dirt, and the overlay's parent
reco2 files are not on local disk. Anything requiring realistic background
therefore still runs on PeLEE.

\subsection{Signal definition used in this section}

\Szero, $\numubar$ charged current, $\Lam \to p\pi^-$ in truth, muon momentum
above $0.15$~GeV/c, and \textbf{both} the reconstructed and the true neutrino
vertex inside the fiducial volume. The branching ratio matters: only
$\sim 60\%$ of \Szero\ events have $\Lam \to p\pi^-$ at all (PDG $63.9\%$), so
the remaining $40\%$ contain no proton or pion to reconstruct and must not be
counted as reconstruction failures.

\paragraph{A fiducial-volume trap.} Requiring only the \emph{reconstructed}
vertex in the FV gives badly wrong answers, because events whose reco vertex is
in the FV but whose true vertex is not are mis-reconstructed --- the reco vertex
belongs to a different interaction --- and they dominate the sample. With reco
FV only the muon appears to be reconstructed $53.7\%$ of the time and all four
particles $10.2\%$; requiring both fiducial cuts gives $91.5\%$ and $18.1\%$.
The internal check that catches this: a ceiling \emph{below} the selection's own
preselection efficiency ($15.6\%$) is impossible.

\subsection{Per-particle reconstruction}

Table~\ref{tab:recoeff} gives, for each of the four final-state particles, how
often a truth-matched PFP exists and where Pandora places it in its hierarchy.
Statistics: 3704 signal events, full NuWro sample.

\begin{table}[h]\centering
\begin{tabular}{lrrrr}
\toprule
particle & primary (gen 2) & secondary (gen $\geq3$) & not reconstructed & total \\
\midrule
$\mu^+$   & $91.5\%$ & $0.3\%$ & $8.2\%$  & $91.8\%$ \\
$p$       & $58.1\%$ & $2.8\%$ & $39.1\%$ & $60.9\%$ \\
$\pi^-$   & $45.7\%$ & $1.6\%$ & $52.8\%$ & $\mathbf{47.3\%}$ \\
$\gamma$  & $59.8\%$ & $2.6\%$ & $37.6\%$ & $62.4\%$ \\
\bottomrule
\end{tabular}
\caption{Reconstruction of the four \Szero\ final-state particles.
The $\pi^-$ is the weak link.}
\label{tab:recoeff}
\end{table}

Found together: $0$ of $4$ in $4.5\%$ of events, $1$ in $7.5\%$, $2$ in
$27.2\%$, $3$ in $42.6\%$, and \textbf{all four in $18.1\%$}. Since the
selection requires a muon, proton, pion and shower candidate, $18.1\%$ is its
hard ceiling; it reports $15.6\%$ at the \texttt{NTracksCut} preselection only
because it fills those slots with whatever PFP is available.

\paragraph{Secondaries do not explain the loss.} The analyser originally
discarded every PFP with generation $\neq 2$. The natural hypothesis was that
the \Lam\ daughters, produced at a displaced vertex, were being reconstructed as
generation-3 secondaries and thrown away. The module was changed to keep them
(labelled by a new \texttt{pfp\_generation} branch, so
\texttt{pfp\_generation}~$==2$ recovers the old behaviour exactly) and the
measurement repeated: secondaries are $6.6\%$ of PFPs but recover only $2.8\%$
for the proton and $1.6\%$ for the pion, and nothing at all for the muon and
photon. The losses are genuine reconstruction failures.

\subsection{Track/shower misclassification}

In events with at least four PFPs (622 events), each truth-matched particle was
compared against the track/shower split the selection uses
(\texttt{pfp\_trk\_shr\_score} $> 0.5$):

\begin{table}[h]\centering
\begin{tabular}{lrrrr}
\toprule
particle & expected & as track & as shower & not found \\
\midrule
$\mu^+$  & track  & $95.7\%$          & $1.6\%$           & $2.7\%$ \\
$p$      & track  & $70.3\%$          & $8.2\%$           & $21.5\%$ \\
$\pi^-$  & track  & $54.2\%$          & $\mathbf{11.6\%}$ & $34.2\%$ \\
$\gamma$ & shower & $\mathbf{15.0\%}$ & $69.5\%$          & $15.6\%$ \\
\bottomrule
\end{tabular}
\caption{Classification of truth-matched particles. Bold entries are
misclassifications that directly corrupt candidate assignment.}
\label{tab:misclass}
\end{table}

All four particles are \emph{found} in $41.2\%$ of these events but all four are
\emph{correctly classified} in only $26.7\%$: the hard split at $0.5$ costs
$14.5$ percentage points, a third of otherwise-perfect events. Two consequences
matter. Photons pass as tracks $15\%$ of the time and then compete for the pion
slot --- which is the origin of the $25\%$ photon contamination measured
independently among pion candidates. And $\pi^-$ are called showers $11.6\%$ of
the time, which removes them from the track pool entirely so they can never be
selected.

\subsection{The pion candidate, and why ranking cannot fix it}

On signal events reaching the pion stage the candidate is a true $\pi^-$ only
$34.2\%$ of the time (muon $89.5\%$, photon $77.6\%$, proton $64.5\%$), with
contaminants $\gamma$ $25.0\%$, $p$ $16.2\%$, $\mu^+$ $10.5\%$. Two hard
constraints make this not a ranking problem:

\begin{itemize}
  \item A true $\pi^-$ PFP exists in only $56.6\%$ of these events (track-like
        in $52.6\%$). That is the ceiling: $43\%$ of the time there is no pion
        to find and the selection assigns one regardless.
  \item $76.3\%$ of events have exactly \emph{one} track-like PFP left after the
        muon and proton are taken, so there is no choice to make. Adding
        wiggliness and $dE/dx$ terms to the ranking changed the outcome by
        exactly zero events across seven weightings.
\end{itemize}

Raising the pion track-score threshold lifts purity from $34.2\%$ to $41.6\%$
but the figure of merit falls monotonically (ratio $1.000 \to 0.442$ over
$0.50 \to 0.95$), so it is not worth doing.

\subsection{The \Lam\ displaced vertex is not usable}

The \Lam\ displaced-vertex handle fails on physics grounds, not implementation.
Measured with truth-matched $p$ and $\pi^-$ tracks ($n=266$):

\begin{itemize}
  \item True \Lam\ flight distance: median $1.85$~cm, mean $3.23$~cm. Not
        $7.89$~cm --- $c\tau$ is the rest-frame value and the \Lam\ from a \Szero\
        decay is slow ($\beta\gamma \approx 0.4$), so $\beta\gamma c\tau \approx
        3$~cm.
  \item Reconstructed primary-vertex resolution: median $5.54$~cm.
  \item Correlation between reconstructed and true flight distance:
        $r = -0.014$.
\end{itemize}

The ruler is three times longer than the quantity being measured. Three
estimators were compared against the true decay vertex --- the midpoint of the
two track starts (what the code used), the distance of closest approach of the
two track lines, and the best of the four endpoint pairings --- and all three
gave a residual of $\sim 5.5$~cm, because all are limited by the same vertex
resolution. The tracks themselves converge well (DCA median $0.33$~cm); it is
the baseline that is poor. This closes the earlier null result
(\texttt{LambdaDecayDist} signal $10.19$~cm against background $11.02$~cm) and
the requirement is withdrawn from the work plan. It would still work for
directly produced, faster \Lam, so the negative result should not be
generalised.

% =============================================================================
\section{A global hypothesis fit}
\label{sec:hypfit}
% =============================================================================

\subsection{Motivation}

The selection currently identifies the muon, then the proton, then the pion,
each conditioned only on what has already been taken, and never uses the fact
that $p\pi^-$ must form a \Lam\ and $\Lam\gamma$ a \Szero. The particles are not
independent and are being treated as if they were. The per-role PID templates in
Sec.~\ref{sec:phase0} show why this fails specifically on the pion: the muon,
proton and photon each have a strong, distinctive signature, while
\textbf{the $\pi^-$ has none} --- it lies between the other three on every
variable measured. It can only be identified by exclusion once the others are
pinned, together with the \Lam\ mass constraint, which is precisely what a
sequential cut chain cannot do.

The proposed replacement is:
\begin{enumerate}
  \item \textbf{Relaxed candidate pool.} Drop \texttt{trk\_score} $>0.5$ as an
        eligibility gate; every PFP is eligible for every role and the score
        becomes a likelihood term. This recovers the $11.6\%$ of $\pi^-$ called
        showers and stops photons monopolising the pion slot.
  \item \textbf{Enumerate assignments} of four roles to $N$ PFPs, $P(N,4)$.
        Computationally free: $N=4 \to 24$, $N=6 \to 360$, $N=8 \to 1680$.
  \item \textbf{Score each hypothesis} with a combined $-2\ln L$ built from the
        two mass constraints, per-role PID templates, muon-at-vertex, the
        $p$--$\pi$ DCA, a length-ordering prior and a soft-photon prior.
  \item \textbf{Partial hypotheses} for the $29.1\%$ of signal with three PFPs,
        where the dominant topology is $2$ tracks $+$ $1$ shower ($51.9\%$),
        i.e.\ one hadron missing with the photon still present.
\end{enumerate}

The headroom is specific: all four particles are found in $41.2\%$ of
$N\geq4$ events but correctly classified in $26.7\%$, so perfect assignment
given what is found is worth a factor $1.54$, before any gain from the enlarged
pool or from partial hypotheses.

\subsection{Phase 0: measured ingredients}
\label{sec:phase0}

Everything the fit is parameterised by has been measured on truth-matched PFPs
and recorded in \texttt{HypFitParams.h} (700 files, 1323 signal events).

\begin{table}[h]\centering
\begin{tabular}{lrrr}
\toprule
constraint & $n$ & Gaussian centre & $\sigma$ \\
\midrule
$m(p,\pi)$ vs \Lam\ ($1.11568$)      & 397 & $1.1234$ & $27.2$~MeV \\
$m(\Lam,\gamma)$ vs \Szero\ ($1.19264$) & 243 & $1.1812$ & $45.4$~MeV \\
\bottomrule
\end{tabular}
\caption{Mass constraints from correctly assigned candidates. The \Szero\ term
is $\sim1.7\times$ wider, as expected from the soft photon.}
\label{tab:massres}
\end{table}

The reconstruction is biased by $+7.7$~MeV on the \Lam\ and $-11.4$~MeV on the
\Szero, which is $28\%$ and $25\%$ of the respective widths. \textbf{The $\chi^2$
must therefore be centred on the measured means, not the PDG masses}, or every
hypothesis carries a built-in offset. Both biases have identifiable causes: the
pion momentum is taken from range under the \emph{muon} hypothesis, and the
photon energy is under-measured (reco $0.0792$ vs true $0.0861$~GeV).

Remaining ingredients:
\begin{itemize}
  \item $p$--$\pi$ distance of closest approach: median $0.35$~cm, $16$--$84\%$
        range $0.08$--$1.36$~cm. Strongly peaked and one of the most powerful
        terms available --- and note that it works precisely where the \Lam\
        decay \emph{distance} fails.
  \item The muon is the longest of the three tracks $92.1\%$ of the time: a soft
        prior worth $\ln(0.921/0.079) = 2.46$ in log-likelihood, not a cut.
  \item $98.9\%$ of PFPs carry \emph{both} a track fit and a shower fit, so the
        relaxed pool requires no re-reconstruction --- only a change of
        eligibility rule.
\end{itemize}

\begin{table}[h]\centering\small
\begin{tabular}{lrrrr}
\toprule
variable & $\mu^+$ & $p$ & $\pi^-$ & $\gamma$ \\
\midrule
\texttt{trk\_shr\_score}    & $1.00$  & $0.90$  & $0.80$  & $\mathbf{0.11}$ \\
\texttt{trk\_llrpid}        & $\mathbf{0.96}$ & $-0.29$ & $0.14$ & $0.56$ \\
\texttt{trk\_pid\_chipr}    & $180$   & $\mathbf{16.4}$ & $96$ & $192$ \\
\texttt{three\_plane\_dedx} & $2.47$  & $\mathbf{10.3}$ & $7.52$ & $3.30$ \\
\texttt{rms\_deflection}    & $0.01$  & $0.01$  & $0.01$  & $\mathbf{0.08}$ \\
\texttt{moliere\_avg}       & $1.78$  & $6.18$  & $8.95$  & $\mathbf{13.3}$ \\
\texttt{trk\_length}        & $\mathbf{139}$ & $12.8$ & $10.8$ & $14.2$ \\
\bottomrule
\end{tabular}
\caption{Per-role PID templates, medians on truth-matched PFPs
($n = 1213/808/650/814$). The muon, proton and photon each have a distinctive
signature; the $\pi^-$ has none.}
\label{tab:templates}
\end{table}

\subsection{The mass-sculpting risk}

Selecting the combination that best fits $m_\Lam$ will sculpt a \Lam-like peak
out of pure combinatorics in background events, manufacturing the signal being
searched for. This is the standard failure mode of combinatoric selection and it
is the main risk in the whole scheme.

\textbf{Decision: the \Lam\ and \Szero\ invariant masses are used for assignment
only and are kept out of the BDT input list.} The best-hypothesis score and the
margin to the second-best hypothesis are permitted as inputs. Before the fit is
adopted, the full machinery must be run on a background-only sample and the
resulting mass distributions inspected for an induced peak. A secondary check is
to assign using one constraint and test with the other.

A second, smaller risk: the enlarged candidate pool raises background acceptance
as well as signal efficiency, so the figure of merit must be re-measured end to
end rather than judged on efficiency alone.

% =============================================================================
\section{Fake data by Poisson throw}
\label{sec:fakedata}
% =============================================================================

\paragraph{The problem.} The current ``data'' sample is the signal file itself
(\texttt{SigmaSelection.C} maps the same NuWro hyperon ntuple to both sample~0 and
sample~4), scaled by POT. That is not a closure test: it contains no background, no
statistical fluctuation, and is perfectly correlated with the signal prediction it
would be compared against.

\paragraph{The fix.} Adopt the same Poisson-throw scheme the CC$1\mu1\pi$ analysis
uses (\texttt{macros/throw\_perrun\_*.C}). For each MC event with central-value
weight $w_{\mathrm{cv}}$ and POT scale $\alpha = \mathrm{POT}_{\mathrm{data}} /
\mathrm{POT}_{\mathrm{MC}}$, draw
\begin{equation}
  n \sim \mathrm{Poisson}\!\left( w_{\mathrm{cv}}\,\alpha \right)
\end{equation}
and write the event $n$ times into the fake-data tree with all weights reset to
unity. Applied across \emph{all} samples -- signal, other hyperons, background,
dirt -- this produces a fake dataset that
\begin{itemize}
  \item contains the correct background admixture, which matters enormously here
        because background outnumbers signal by $\sim\!900{:}1$;
  \item carries the right statistical fluctuations, so uncertainties can be
        validated rather than assumed;
  \item is statistically independent of the prediction, making a genuine closure
        test possible.
\end{itemize}

\paragraph{Implementation.} A \texttt{throw\_sigma0.C} modelled on
\texttt{throw\_perrun\_w.C}: chain the processed MC, drop the multisim weight
branches, \texttt{CloneTree(0)} to preserve the \Szero\ observables, throw, and
stamp \texttt{summed\_pot} with the target POT. One seed per sample so the throw is
reproducible.

\paragraph{Caveat to keep in view.} A Poisson throw from the same MC is a
\emph{statistical} closure test, not a model test. It cannot reveal generator
mismodelling, because the ``data'' and the prediction come from the same generator.
A genuine model test needs a hyperon sample from a different generator -- GENIE
against the current NuWro -- which does not exist locally. Until it does, the
throw validates the statistical and normalisation chain only, and that limitation
should be stated wherever the closure is quoted.

% =============================================================================
\section{Exposure: RHC now, combined later}
\label{sec:exposure}
% =============================================================================

\paragraph{Now: RHC only.} Everything is scaled to the full RHC data exposure,
$\mathrm{POT}_{\mathrm{data}} = 1.1526\times10^{21}$:

\begin{center}
\begin{tabular}{llr}
\toprule
Sample & Source & Generated POT \\
\midrule
\Szero\ signal / hyperons & NuWro hyperon-enhanced & $2.1764\times10^{23}$ \\
Background   & RHC overlay, Runs 1$+$2$+$4a$+$4b$+$4c & $9.94482\times10^{21}$ \\
Dirt         & Run~1 RHC only & $4.21274\times10^{20}$ \\
\bottomrule
\end{tabular}
\end{center}

\noindent Two known imperfections: the dirt sample is Run-1 RHC only, scaled up to
the full RHC exposure, and the signal sample is generated with no cosmic overlay
(NoCR), which is why its non-muon truth matching cannot be trusted.

\paragraph{Later: combined FHC$+$RHC.} The signal is $\numubar$ charged-current,
so RHC is the sensitive mode -- but FHC carries a substantial wrong-sign
$\numubar$ component ($34.0\%$ of the FHC flux, against $53.3\%$ for RHC), and the
FHC exposure is $8.857\times10^{20}$~POT. Folding FHC in therefore adds real signal
rather than only background, and the CC$1\mu1\pi$ analysis already has the
machinery: per-run POT scaling, a combined \texttt{file\_properties}, and combined
flux normalisation.

Deferred deliberately. The combination is a normalisation exercise, and at $0.1\%$
purity it would add background faster than signal in absolute terms. The selection
work of Sec.~\ref{sec:purity} comes first; the exposure extension is worth doing
once the selection is worth extending. When it happens it needs: FHC overlay
samples processed through the same selection, an FHC hyperon signal sample (or a
justified reuse of the RHC one, given the signal is beam-mode independent at the
interaction level), FHC dirt, and a combined-mode flux normalisation following
\texttt{Flux} in the CC$\pi$ xsec configs.

% =============================================================================
\section{Work plan}
% =============================================================================

\begin{enumerate}
  \item \textbf{Selection (priority): global hypothesis fit.} Relaxed candidate
        pool, combinatoric role assignment scored against the \Lam\ and \Szero\
        mass constraints and the per-role PID templates, then partial hypotheses
        for three-PFP events. Sec.~\ref{sec:hypfit}. Phase~0 (resolutions,
        templates, priors) is complete and recorded in
        \texttt{HypFitParams.h}; Sec.~\ref{sec:phase0}.
        \emph{The \Lam\ displaced-vertex requirement listed in earlier versions
        of this plan is withdrawn: the true flight distance (median $1.85$~cm) is
        three times smaller than the vertex resolution ($5.54$~cm) and the
        reconstructed quantity is uncorrelated with truth, $r=-0.014$.
        Sec.~\ref{sec:recostudy}.}
  \item \textbf{Framework port.} Branch bindings, \texttt{Sigma0} selection class,
        factory registration, cut-flow validation against the standalone code
        before any config work. Sec.~\ref{sec:port}.
  \item \textbf{Background through HyperonProduction.} A hypana-only analysis is
        blocked: hypana output exists for the hyperon samples only, and the RHC
        overlay's parent reco2 files are not on local disk. Until they are
        retrieved and processed, anything needing realistic background runs on
        PeLEE. Sec.~\ref{sec:recostudy}.
  \item \textbf{Fake data.} \texttt{throw\_sigma0.C}, all samples, RHC POT.
        Sec.~\ref{sec:fakedata}.
  \item \textbf{Systematics.} Single-bin \texttt{univmake} for the background
        covariance.
  \item \textbf{Combined exposure.} FHC$+$RHC once the selection justifies it.
        Sec.~\ref{sec:exposure}.
  \item \textbf{Limit.} Outside the framework, reading its covariance.
\end{enumerate}

\end{document}
